% Journal-style table fragments for the Trust, Autonomy, and Evidence manuscript.
% Requires: \usepackage{booktabs,tabularx,threeparttable}

\begin{table}[t]
\centering
\begin{threeparttable}
\caption{Derivation of the event-level practical-control chain}
\label{tab:construct-derivation}
\footnotesize
\begin{tabularx}{\linewidth}{l X X X}
\toprule
Stage & Why required & Minimum record & Error prevented \\
\midrule
Information access & Judgment cannot use information the person never received. & Timestamped delivery, interface, or access record. & Treating post-action notice as review. \\
Comprehension capacity & Delivery does not establish capacity to interpret. & Training, explanation, test, or reasoning record. & Treating visibility as understanding. \\
Intervention authority & Control requires a right to change the decision. & Delegation, procedure, or permission record. & Treating responsibility as permission. \\
Intervention feasibility & A formal right may be unusable under actual conditions. & Timing, staffing, latency, access, or workload record. & Treating theoretical authority as usable authority. \\
Exercised judgment & Available control does not establish that judgment occurred. & Approve, reject, modify, stop, or escalate record. & Inferring action from a role assignment. \\
Execution propagation & A judgment that never reaches execution cannot control the event. & Linked system or institutional execution record. & Inferring operational change from an unexecuted judgment. \\
\bottomrule
\end{tabularx}
\begin{tablenotes}[flushleft]\footnotesize
\item The released assessment JSON uses \texttt{effect} for execution propagation. The manuscript term avoids implying that the change was beneficial or caused the final outcome.
\end{tablenotes}
\end{threeparttable}
\end{table}

\begin{table}[t]
\centering
\begin{threeparttable}
\caption{Assessment-state decision rules}
\label{tab:state-rules}
\begin{tabularx}{\linewidth}{l c X}
\toprule
State & Code & Decision rule \\
\midrule
Supported & S & Direct, contemporaneous evidence satisfies the stage definition. \\
Partially supported & P & Some required elements are present and a material gap remains. \\
Unsupported & U & Available evidence contradicts the stage or shows that the condition was absent. \\
Indeterminate & I & The packet lacks enough evidence to decide. \\
Outside scope & O & The stage does not apply within the declared case boundary, with a written reason. \\
\bottomrule
\end{tabularx}
\begin{tablenotes}[flushleft]\footnotesize
\item A missing record produces an indeterminate state unless the protocol or case design establishes that the record should exist. The states are categorical. No numeric distance or aggregate score is assigned.
\end{tablenotes}
\end{threeparttable}
\end{table}

\begin{table}[t]
\centering
\begin{threeparttable}
\caption{Practical-control states across three public cases}
\label{tab:practical-control}
\begin{tabular}{l c c c}
\toprule
Stage & Oko, 1983 & ZG710, 2003 & F/A-18C, 2003 \\
\midrule
Information access & P & P & P \\
Comprehension capacity & P & U & I \\
Intervention authority & P & S & S \\
Intervention feasibility & P & U & I \\
Exercised judgment & P & U & I \\
Execution propagation & P & U & U \\
\midrule
Correction & O & O & O \\
Repair & O & U & U \\
Institutional reform & P & S & S \\
\bottomrule
\end{tabular}
\begin{tablenotes}[flushleft]\footnotesize
\item S = supported; P = partially supported; U = unsupported; I = indeterminate; O = outside scope. The first six rows determine the event-control result. Correction, repair, and institutional reform describe post-event response and do not alter that result. The table reports 27 item-level findings from three purposefully selected cases and supplies no frequency, causal, or population estimate.
\end{tablenotes}
\end{threeparttable}
\end{table}

\begin{table}[t]
\centering
\begin{threeparttable}
\caption{Formal authority compared with event-level practical control}
\label{tab:authority-event-control}
\begin{tabularx}{\linewidth}{l c c X}
\toprule
Case & Authority & Event control & Reason \\
\midrule
Oko, 1983 & P & Unresolved & All six required stages are partially supported. \\
Patriot ZG710, 2003 & S & Fail & Four required stages are unsupported. \\
Patriot F/A-18C, 2003 & S & Fail & Execution propagation is unsupported; three stages are indeterminate. \\
\bottomrule
\end{tabularx}
\begin{tablenotes}[flushleft]\footnotesize
\item Supported authority establishes one required condition and cannot substitute for the other five. Results are derived from released states under the v0.16.0 rule.
\end{tablenotes}
\end{threeparttable}
\end{table}

\begin{table}[t]
\centering
\begin{threeparttable}
\caption{Formal search and final screening state}
\label{tab:formal-search}
\begin{tabularx}{\linewidth}{l X r l}
\toprule
Stage & Record class & Count & Status \\
\midrule
Retrieval & Direct queries & 184 & Declared open-index queries complete \\
Retrieval & Citation chains & 2,482 & Fourteen of fifteen seeds resolved \\
Pooling & Combined records & 2,666 & Before deduplication \\
Pooling & Deduplicated records & 2,431 & Unit for screening \\
\midrule
Final screening & Retain close & 27 & Confirmed close set \\
Final screening & Retain background & 45 & Thirteen prior plus 32 confirmed \\
Final screening & Exclude single component & 10 & Relevant component only \\
Final screening & Exclude topic & 1,259 & Outside the review question \\
Final screening & Inaccessible & 1,087 & Substantive screening unresolved \\
Final screening & Outside cutoff & 3 & Published after cutoff \\
\midrule
Author gate & Completed queue & 89 & All decisions recorded \\
\bottomrule
\end{tabularx}
\begin{tablenotes}[flushleft]\footnotesize
\item The six final screening classes sum to 2,431. Mark Julius Banasihan is the decision owner for the 89-record queue, with disclosed AI assistance. The 1,087 inaccessible records and authenticated database searching remain separate coverage limits.
\end{tablenotes}
\end{threeparttable}
\end{table}

\begin{table}[t]
\centering
\begin{threeparttable}
\caption{Proposal-to-author decision changes}
\label{tab:author-screening-changes}
\begin{tabular}{l r r r r r}
\toprule
Proposed class & Records & Close & Background & Topic & Single component \\
\midrule
Retain close & 12 & 12 & 0 & 0 & 0 \\
Author attention & 77 & 15 & 32 & 20 & 10 \\
\midrule
Total author queue & 89 & 27 & 32 & 20 & 10 \\
\bottomrule
\end{tabular}
\begin{tablenotes}[flushleft]\footnotesize
\item The author confirmed all 12 proposed close records. The 77 attention records produced 15 additional close sources and 32 background sources. These decisions close the declared author gate. They do not resolve the inaccessible-record or authenticated-database gates.
\end{tablenotes}
\end{threeparttable}
\end{table}

\begin{table}[t]
\centering
\begin{threeparttable}
\caption{Residual-risk retrieval and screening checkpoint}
\label{tab:residual-risk-checkpoint}
\begin{tabularx}{\linewidth}{l X r X}
\toprule
Scope & State & Count & Claim boundary \\
\midrule
Forward citations & Retrieval outcomes recorded & 102 & Frozen stratum complete \\
Forward citations & Recovered content & 71 & Eligible for screening \\
Forward citations & Screening decisions & 71 & 13 close, 22 background, 11 single-component, 25 topic \\
Forward citations & Proposition permissions & 5 & One bounded proposition each \\
Forward citations & Background-only or quarantined & 8 & Two background; six quarantined \\
Direct queries & Screening decisions & 5 of 5 & RS-DQ-004 has zero proposition permission \\
Recovery population & Retrieval outcomes & 107 of 1,087 & 980 outcomes open \\
Recovery population & Recovered-content decisions & 76 of 76 & Screening complete \\
\bottomrule
\end{tabularx}
\begin{tablenotes}[flushleft]\footnotesize
\item These counts describe the frozen recovery workflow and the v0.14 proposition-review overlay. They supply no close-source prevalence estimate for the 1,087-record population because 980 retrieval outcomes remain open. Proposition permission is narrower than source inclusion and does not establish generalizability.
\end{tablenotes}
\end{threeparttable}
\end{table}

\begin{table}[t]
\centering
\begin{threeparttable}
\caption{Versioned correction of the Oko assessment}
\label{tab:oko-correction}
\begin{tabularx}{\linewidth}{l c c X}
\toprule
Stage & v0.3.0 & v0.6.0 & Material gap recorded in v0.6.0 \\
\midrule
Information access & S & P & No contemporaneous delivery, interface, or command record was located. \\
Comprehension capacity & S & P & No contemporaneous reasoning, review, or explanation record was located. \\
Intervention authority & S & P & No contemporaneous delegation or command-procedure record was located. \\
Intervention feasibility & S & P & No contemporaneous timing or operating record was located. \\
Exercised judgment & S & P & No contemporaneous decision or communication log was located. \\
Execution propagation & S & P & No contemporaneous linked action, stop, or escalation record was located. \\
\bottomrule
\end{tabularx}
\begin{tablenotes}[flushleft]\footnotesize
\item The packet remained fixed. The v0.6.0 adjudication applied a protocol frozen before reassessment and admitted no new historical source. The correction records a change in evidence classification.
\end{tablenotes}
\end{threeparttable}
\end{table}

\begin{table}[t]
\centering
\begin{threeparttable}
\caption{Availability of coding-stability evidence}
\label{tab:coding-stability}
\begin{tabularx}{\linewidth}{X c c c}
\toprule
Test condition & Oko, 1983 & ZG710, 2003 & F/A-18C, 2003 \\
\midrule
A second coding exists & Yes & No & No \\
The same source packet was used & Yes & NA & NA \\
The same evidence rule was used & No & NA & NA \\
Comparable event-level stages & 6 & 0 & 0 \\
Unchanged classifications & 0 & NA & NA \\
Changed classifications & 6 & NA & NA \\
Independent second assessor & No & No & No \\
Reliability claim eligible & No & No & No \\
\bottomrule
\end{tabularx}
\begin{tablenotes}[flushleft]\footnotesize
\item Oko's six changes arose when the direct-and-contemporaneous rule was applied to the frozen packet. The comparison records a correction under a changed classification rule. It cannot estimate intra-rater stability or inter-rater reliability. The two Patriot cases have one released coding each.
\end{tablenotes}
\end{threeparttable}
\end{table}
